\usepackage[utf8]{inputenc}
\usepackage{lmodern}
\usepackage{amsmath}
\usepackage[ngerman]{babel}
\usepackage[ngerman]{translator} % Weitere Übersetzungen
\usepackage{graphicx}
\usepackage{layout}
\usepackage{fancyhdr}
\usepackage{roboto}
\usepackage{sectsty}
\usepackage{lastpage}
\usepackage{tabularx}
\usepackage[table]{xcolor}
\usepackage{tikz}
\usepackage{array}
\usepackage{lipsum}
\usepackage{multirow}
\usepackage{float}
\usepackage{caption}
\usepackage{hyperref}
\usepackage{geometry}            % Dient zur Interface dimensionierung
\usepackage{pdflscape}           % bessere Lesbarkeit des PDF Dokuments
\usepackage{typearea}  % loaded automatically if using KOMA-Script class

\usepackage{pgfgantt}

\hypersetup{%
    colorlinks=false,% hyperlinks will be black
    pdfborderstyle={/S/U/W 0}% border style will be underline of width 1pt
}

%Tabellen Caption links orientiert
\captionsetup[table]{skip=0pt,singlelinecheck=off}

%Tabellen- & Abbildungsverzeichnis als Section und in Inhaltsverzeichnis 
\makeatletter
\renewcommand\listoftables{%
    \section{\listtablename}%
    \@mkboth{\MakeUppercase\listtablename}%
        {\MakeUppercase\listtablename}%
    \@starttoc{lot}%
}
\renewcommand\listoffigures{%
    \section{\listfigurename}%
    \@mkboth{\MakeUppercase\listfigurename}%
        {\MakeUppercase\listfigurename}%
    \@starttoc{lof}%
}
\makeatother

%Weitere Spalten typen
\newcolumntype{L}[1]{>{\raggedright\let\newline\\\arraybackslash\hspace{0pt}}m{#1}}
\newcolumntype{C}[1]{>{\centering\let\newline\\\arraybackslash\hspace{0pt}}m{#1}}
\newcolumntype{R}[1]{>{\raggedleft\let\newline\\\arraybackslash\hspace{0pt}}m{#1}}

%Seitenformatierung anpassen
\textwidth=400pt
\textheight=650pt
\oddsidemargin=25pt
\footskip=20pt
\fancypagestyle{plain}{}
\pagestyle{fancy}
\voffset=0pt

%Sans seriff Font
\allsectionsfont{\sffamily}

%Header höhe anpassen
\setlength\headheight{33pt}
%Bilderpfad
\graphicspath{{bilder/}}
%Header und Footer
\fancyhf{}
\cfoot{\today}
\rfoot{Seite \thepage \hspace{1pt} von \pageref{LastPage}}
\lfoot{Version 0.0.1}
\chead{Titel der IPA}
\lhead{\includegraphics[width=3cm]{puzzle_tagline_blau}}  
\rhead{Sylvain Gilgen}

%Weiteres
\definecolor{puzzleblue}{RGB}{30,90,150}
\author{Sylvain}
\title{IPA}